\section{Two Derivations}
\label{sec:derivations}

From the core definition we derive two structural results. Both concern the
\emph{choice of basis} in which the round trip is expressed.

\subsection{Derivation 1: Basis Choice (why perturbation and hierarchy are equivalent)}
\label{sec:basis}

Perturbation series and hierarchical (level-by-level) structures are both
``leading term plus successive corrections''---the same series structure,
algebraically equivalent. They differ only in the \emph{semantics of the order}:
precision in perturbation, concept in hierarchy. When conceptual crossings
(coupling terms) exist, the equivalence breaks down and a joint basis is
required.

\paragraph{The basis is freely choosable.}
Among perturbation bases, hierarchical bases, tree bases, and Hilbert-space
bases, the optimal basis is the one in which the expansion coefficients are most
sparse (least redundant)---i.e., the most compressed representation. Sparsity of
coefficients is thus the criterion for basis optimality, and it is exactly the
round-trip condition applied to the representation itself.

\paragraph{Hierarchy/tree corresponds strictly to multiresolution bases.}
The rigorous counterpart of hierarchy and tree is multiresolution analysis
(Haar, Daubechies, Faber--Schauder)~\cite{mra_daubechies}: compared with
Fourier bases, multiresolution bases add \emph{locality} (basis functions with
local support) and \emph{nestedness} ($V_j \subset V_{j+1}$, tier $=$ scale).
``Jumping levels'' corresponds to exponential scales $2^{j}$ combined with
tree-like branching. Orthogonality can be preserved (Haar and Daubechies bases
are $L^{2}$-orthogonal), so the round-trip identity is not sacrificed.

\paragraph{Multiplicative structure as tensor products.}
Multiplicative (factorized) structure unifies with the same picture: a single
rooted tree is branching; a product of multiple roots is a tensor product
($1\mathrm{D}\otimes1\mathrm{D}\otimes\cdots$), and the two are unified under
multiresolution analysis. This answers how to describe a factorized structure
that is not a single-root split: as a tensor product of rooted trees, one
factor per root.

\subsection{Derivation 2: Two-Stage Compression (domain selection then refinement; order cannot be reversed)}
\label{sec:twostage}

General (structural) compression is responsible for \emph{selecting the domain};
formula-level compression is responsible for \emph{refining within the domain}
(lossless derivation). The ``losslessness'' of formula-level work is relative to
the domain of applicability, and that domain must be fixed by general thinking.
Reversing the order---refining before selecting the domain---means deriving
precisely in the wrong domain, which drives the result farther off, not closer.

This explains an empirical observation: in some respects formula-level methods
are better while general thinking is not. Formula-level compression excels in
the refinement stage (precise, derivable) but helps little---and can even harm---
in the domain-selection stage, by inducing premature priors.

\paragraph{Domain selection is a three-step process.}
Domain selection is not a static judgment but a process: (i) verify the range in
which the axioms hold; (ii) transfer the axioms that still hold, together with
their conditions; (iii) think about and experiment with new conditions of
applicability. This fills the gap in the reasoning chain between induction and
deduction (induction $\to$ condition verification $\to$ transfer $\to$
experiment $\to$ deduction $\to$ predictive quality), where previously only
deduction was explicit.

\paragraph{Meta-criterion.}
Any mathematical metaphor is itself ``compression plus selection''. Its value is
measured by whether it yields falsifiable predictions: if yes, it is useful; if
no, it is empty compression and should be discarded.
